% page 488 of the facsimile (image p-487 of the scan)
% block 165.1 x 259.3 mm ; left margin 8.0 mm
\pgc{-12.700mm}{0.000mm}{
\l{}
\l{}
\l{\cel{3}480}
\l{}
\l{radiko.\cel{2}I. (botLLa parto infra di vejetanto, ordinare sinkanta}
\l{\cel{3}aden tero, ube lu su developas inversa-sinse di la stipo,}
\l{\cel{3}e di qua la rolo esas fixigar la planto en sulo e pumpar\cel{2}la}
\l{\cel{3}elementi de qui nutresas la planto. - Parto di organo (dento}
\l{\cel{3}haro, pilo, unglo) per qua lu su enplantacas en altra parto.}
\l{\cel{3}II. (\sou{metaf.)}\cel{2}(a) (\sou{gram}.) Vorto primitiva, de qua derivesas}
\l{\cel{3}altra vorti. (b)(\sou{matem.})\cel{1}\sou{Radiki di equaciono}\cel{1}: quanti quin}
\l{\cel{3}onu povas substitucar a la nekonocaji, kom satisfacanta ega-}
\l{\cel{3}le la donaji di la equaciono. -\cel{1}\sou{Extraktar radiko, quadrata,}}
\l{\cel{3}\sou{kubala, di nombro}\cel{1}: Trovar la nombro qua, multiplikite unople}
\l{\cel{3}o duople ye su ipsa, donis ta nombro.\cel{2}EFIS.}
\l{}
\l{\sou{radiografar}. (\sou{trans.)}\cel{2}Fotografar la internajo di la homo-korpo}
\l{\cel{3}per la X-radii (Roentgen-radii) - por skopi medicinala, kirur-}
\l{\cel{3}giala, e c. - EFIRS.}
\l{}
\l{\sou{radiometro}. Aparato uzata por mezurar la intenseso-grado di la}
\l{\cel{3}lumo-radii. - EFIS.}
\l{}
\l{\sou{radioskopar}. (\sou{trans.})\cel{2}(\sou{medicino})\cel{2}Explorar la ombri quin produk-}
\l{\cel{3}tas la X-radii (Roentgen-radii) sur skrino fluorecanta ek}
\l{\cel{3}platino-cianuro di bario o di tungsteno. - EFIS.}
\l{}
\l{\sou{radirar}. (\sou{trans.)}\cel{3}Grabar, facante la desegnuro per stal-instru-}
\l{\cel{3}mento stipeto-forma, kun punto tre akuta por povar trasar sur}
\l{\cel{3}la varniso o la vaxo-strato linei tre dina e delikata. - D.}
\l{}
\l{\sou{radiumo.}\cel{2}(\sou{kem.})\cel{2}Elemento metala, tre skarsa, deskovrita en}
\l{\cel{3}1898 da Pierre e Marie Curie kun J. Bémont, en la rezidui ba-}
\l{\cel{3}rioza obtenita lor la trakto di pekblendo (oxido naturala di}
\l{\cel{3}uranio). -\cel{1}\sou{Simbolo kemiala}\cel{1}: Ra.\cel{3}- DEFIRS.}
\l{}
\l{\sou{radiuso.}\cel{2}(\sou{anat.)}\cel{1}Ta, ek la du osti di la avanbrakio, qua esas}
\l{\cel{3}sur la latero vers-extere. - EFS.}
\l{}
\l{\sou{radotar}. (\sou{netrans., pri}). Parolar nekonsequante, nekoherante,}
\l{\cel{3}konseque de feblesko di la spirito. - EF.}
\l{}
\l{\sou{rafano}. (\sou{bot.)}\cel{3}Speco di kruciferi, di la genero \textquotedbl{}cochlearia\textquotedbl{}:}
\l{\cel{3}planto perena, kun grosa radiko pulpoza, blanka, qua (freshe}
\l{\cel{3}raspita) povas remplasar mustardo. - L.\cel{1}\sou{raphanus sativus}.- I}
\l{}
\l{\sou{rafido}. (\sou{biol.)}\cel{3}Kristali aguloforma quin onu observas en ula}
\l{\cel{3}celuli di la animali o di la vejetanti.}
\l{}
\l{\sou{rafinar}. (\sou{trans.)}\cel{2}I. Igar plu pura per deprenar la substanci}
\l{\cel{3}stranjera quin (lu) kontenas. - II. (\sou{metaf.}) Igar plu delika-}
\l{\cel{3}ta, plu subtila. - DEFIRS.}
\l{}
\l{\sou{raflo.}\cel{2}(\sou{bot.)}\cel{2}Vito-grapo de qua onu deprenis lua beri. - F.}
\l{}
\l{\sou{rafto.}\cel{2}Asemblitaro de trabi (o de arbori abatita) interligita}
\l{\cel{3}tale ke li formacas quaza planko-sulo flotaciva. - E.}
\l{}
\l{\sou{rago}.\cel{2}Fragmento de stofo, telo, lano, e c., qua pendas,mi-lace}
\l{\cel{3}rite, de olda vesto par-konsumita. -\cel{2}E.}
}
